% Original attached proof of night:P3F.3; retained solely as audit evidence.
\begin{proof}
\begin{enumerate}[label=\arabic*.]
\item The primary source is Huize Jin, arXiv:2608.26698v2 (12September2026), Theorem1.2, Corollary1.3, Proposition3.1, Lemmas4.4 and4.8, and Theorem4.9. The downloaded PDF is retained in \path{records/primary_sources}. The ambient manifold is the oriented product \(M=S^2\times S^3\). Its orientation-preserving Torelli group is

\[
 I(M)=\mathbb Za\oplus\mathbb Zb.
\]

\item Here \(a=t_{S_1,\alpha}\) twists a standard \(S^2\times D^3\), and \(b=t_{S_1\#S_1,\alpha}\) twists the connected sum of two oriented copies of that core \(S^2\). The normalization of \(\alpha\in\pi_3SO(3)\) is that the rank3 bundle over (S\^{}4) clutched by \(\alpha\) has first Pontryagin number4. The twist is explicitly

\[
 (x,y)\longmapsto(\alpha(y)x,y),
\]

where \(\alpha:D^3\to SO(3)\) equals the identity near its boundary. Extending by the identity gives the same formula with \(\alpha:S^3\to SO(3)\). In particular the core sphere has dimension2.

\item For \(f\in I(M)\), Jin uses

\[
 M_f=M\times[0,1]/((x,0)\sim(f(x),1)),
\]

with the preferred spin structure specified by a fixed disc and a class \(x\in H^2(M_f;\mathbb Z)\) restricting to the positive generator on \(S^2\). Put

\[
 k(f)=\langle x^3,[M_f]\rangle,\qquad
 \ell(f)=\langle p_1(TM_f)x,[M_f]\rangle,
 \qquad A(f)=\frac{4k(f)-\ell(f)}{24}.
\]

\item The theorem identifies \(f\mapsto(A(f),k(f))\) with an isomorphism \(I(M)\to\mathbb Z^2\). Its normalization gives

\[
 (k(a),\ell(a))=(1,4),\qquad
 (k(b),\ell(b))=(8,8),
\]

so \(b-8a\) and \(a\) map to the standard ordered basis of \(\mathbb Z^2\). In particular

\[
 [f]=(k(f)-8A(f))a+A(f)b. \tag{P3F.3.1}
\]

\item These assertions follow from the cited theorem by direct substitution, and verify that the congruence is \(\ell\equiv4k\pmod{24}\). The pair ((-1,-4)) gives ({[}f{]}=-a). It is a nonzero element of \(I(S^2\times S^3)\); its effect on the associated glued 7-manifold is computed explicitly next.

\end{enumerate}
\end{proof}

% Original attached proof of night:P3F.7; retained solely as audit evidence.
\begin{proof}
\begin{enumerate}[label=\arabic*.]
\item For the AC26 sphere member, its delta circle is semifree, has one fixed curve F=P1, and its orbit pair is diffeomorphic to the standard pair \(S5,S2\). Its fixed normal bundle is O\(-1\) plus O\(-1\) with respective weights\(-1,+1\). Replacing that fixed neighbourhood by the conifold smoothing gives a free circle manifold S3 times S3, with quotient S2 times S3.

\item On an ordinary fibre the delta circle is primitive and acts freely. At a finite cap, a nontrivial deck power has nonintegral gamma character, since gamma\(v_j\) is prime to m\_j. It therefore cannot agree with a delta circle rotation. This proves freeness there. In the cusp fan the directions are ±e1,±e2,±\(e1-e2\). A finite circle stabilizer at an edge of direction d has order dividing det\(d,e2\), which is±1 off the e2 edges; on an e2 edge the entire circle fixes it. At the two-dimensional torus orbit corresponding to a vertex of the height-one triangulation there is no stabilizer. The zero-dimensional torus orbits corresponding to triangles are fixed, and each lies on one of the fixed edges. The period lattice acts transitively on the e2 edges, with no stabilizer, so their quotient is one fixed P1.

\item For the vertical wall use the four rays

\begin{lstlisting}
r0=(0,0,1), r1=(0,1,1), r+=(1,0,1), r-=(-1,1,1).
\end{lstlisting}

\item They satisfy r++r-=r0+r1. The two adjacent unimodular charts have coordinates\(x0,x1,z\) and(x0',x1',z'), with exact transition

\begin{lstlisting}
x0'=z x0,  x1'=z x1,  z'=1/z,
q=x0 x1 z=x0' x1' z'.
\end{lstlisting}

\item The fixed curve is x0=x1=0. The transitions give O\(-1\)\textsuperscript{2}, and the cocharacter e2 has weights\(-1,+1,0\). These identities are independently checked from the ray matrices in code/part3/orbit\_mv.py. The map

\begin{lstlisting}
A=x0, B=z x0, C=x1, D=z x1
\end{lstlisting}

contracts the curve to AD=BC. It intertwines the circle with weights \(-1,-1,+1,+1\) on\(A,B,C,D\). On AD-BC=epsilon nonzero the action is free. For epsilon=1, polar decomposition of SL2\(C\) as U H gives the explicit equivariant diffeomorphism SU2 times R3, with the circle acting on U by left multiplication and acting trivially on H. A compact Milnor neighbourhood is thus S3 times D3 and its quotient is S2 times D3.

\item For completeness, the orbit is identified integrally. Let A=X minus F and B=A/S1. General position gives pi1\(A\)=0. Alexander duality gives H*\(A;Z\)=Z in degrees0,3 and zero otherwise. The oriented normal link S3 is a primitive generator of H3\(A\). Its quotient is the normal link S2, and the circle bundle on these links is the Hopf bundle. In the Gysin sequence, e generates H2\(B\), and integration H3\(A\)-\textgreater H2\(B\) is multiplication by1: evaluate on this linked S2 and its Hopf preimage. Every other positive-degree group of B vanishes. Also pi1\(B\)=0 by the homotopy sequence of the principal circle bundle.

\item Conjugating the negative-weight line converts the fixed normal bundle to a positive-weight complex bundle O\(1\) plus O\(-1\). Its clutching loop diag\(z,z^-1\) in SU2 bounds the explicit disc with SU2 coordinates \((a,b)=(\sin(\pi r/2)e^{i\theta},\cos(\pi r/2))\); joining its centre to the identity gives a nullhomotopy. Thus it is a trivial unitary rank2 bundle. Its orbit normal bundle is a trivial oriented rank3 bundle. With this trivialization the boundary is F times S2, and the Euler class is just the meridian generator, with no F term.

\item Now Y=X/S1 is the union of B (truncated at that boundary) and F times D3. The maps on H0 and H2 in Mayer--Vietoris, with domain H2 ordered (F,meridian) and target(complement,tube), are respectively

\begin{lstlisting}
[1;-1],       [[0,1],[-1,0]].
\end{lstlisting}

\item The latter has determinant1. H4(boundary)=Z maps to zero and gives H5\(Y\)=Z; every other positive intermediate group is zero. All matrices and their Smith forms are printed in orbit\_mv.out. Van Kampen gives pi1\(Y\)=0. Thus Y is a homotopy5-sphere. Kervaire--Milnor1963, the order table on p.~504 and the Smale comparison on p.~505, gives Theta5=0, so Y is diffeomorphic to S5. Haefliger1966 Corollary6.6 says C\_n\^{}q=0 for n\textless2q-3; with n=2,q=3, and Theorem1.2 converting concordance to isotopy, every embedding S2 in S5 is unknotted. This identifies the orbit pair.

\item One may now construct the equivariant identification with the standard circle action on the sphere boundary of D3 times D4. On the free complement the principal circle bundle over S2 times D3 has Euler generator and is the product Hopf bundle S3 times D3. On the fixed tube it is the trivial positive-weight rank2 bundle just constructed. Their two boundary identifications differ by a gauge map from S2 times S2 to S1. Such a map has a global logarithm because H1(S2 times S2;Z)=0, so it extends over a collar and removes the discrepancy. This gives an equivariant diffeomorphism. It is also the specialization of Levine's classification as stated in Harvey--Kerin--Shankar2020 Theorem4.1, whose hypotheses are checked here: smooth codimension4 fixed set and H2\(F;Z\) without2-torsion. Replacing the fixed tube S2 times D4 by the free piece D3 times S3 doubles the free complement, giving S3 timesS3; the quotient doubles S2 timesD3, giving S2 timesS3. The normal framing is unique up to homotopy because pi2\(SO4\)=0. This is the asserted conifold surgery.

\end{enumerate}
\end{proof}

% Original attached proof of night:P3F.22; retained solely as audit evidence.
\begin{proof}
\begin{enumerate}[label=\arabic*.]
\item Use the oriented Hopf bundle with Euler class the positive generator of the oriented two-sphere. Orient its total three-sphere by the circle followed by the base, and orient the four-ball which it bounds accordingly. Use the product orientation on each four-ball and the outward-normal-first convention on \(\partial(D^4_u\times D^4_y)\). For the gluing convention \[
 \lambda(f)=(D^4_u\times S^3_y)\cup_{\widetilde f}
                  (S^3_u\times D^4_y),
 \qquad x\in\partial(D^4_u\times S^3_y)
       \sim\widetilde f(x),
\] where \(\widetilde f\) is a principal-circle lift of \(f\in I(S^2\times S^3)\), the oriented Eells--Kuiper invariant is \[
 \mu(\lambda(f))=\frac{A(f)}{28}\pmod{\mathbb Z},
 \qquad A(f)=\frac{4k(f)-\ell(f)}{24}.                 \tag{P3F.22.1}
\] Consequently \(\lambda\) is a homomorphism to \(\Theta _7\), \(\lambda(a)=0\), and \(\lambda(b)\) is the generator with \(\mu=1/28\). Its kernel is \(\mathbb Za\oplus28\mathbb Zb\).

\item The assertion is a transfer computation for the specified gluing map. An assignment of its value to a parameter loop also requires that loop's actual oriented return map. In particular, this proposition by itself does not select the sign of the OS47 parameter loop.

\subsubsection{1. The marked double and its bordism}\label{the-marked-double-and-its-bordism}

\item Put \(K=S^2\times D^4\), with boundary \(M=S^2\times S^3\), and form \[
 N_f=K_+\cup_f(-K_-),\qquad x\in\partial K_+\sim f(x)\in\partial K_-.
\] Let \(x\in H^2(N_f;\mathbb Z)\) restrict to the positive generator on both copies of \(S^2\). Mayer--Vietoris and van Kampen give \[
 \pi_1(N_f)=0,\qquad H^j(N_f;\mathbb Z)=
 \begin{cases}\mathbb Z&j=0,2,4,6,\\0&\text{otherwise}.
 \end{cases}
\] The two spin structures on the pieces agree on the simply connected boundary, so \(N_f\) is spin. There is a unique \(y\in H^4(N_f;\mathbb Z)\) with \(\langle xy,[N_f]\rangle=1\).

\item For clarity the required bordism is an actual collar construction, the one in Jin26v2 Lemma6.2. Write an inward collar of \(K\) as \(M\times[0,1]_s\), and identify the two collars at the ends of \(K\times[0,1]_t\) by \((z,s,0)\sim(f(z),s,1)\). Rounding corners gives a seven-manifold. Its \(s=0\) boundary is the mapping torus in the convention \((z,0)\sim(f(z),1)\); its other boundary is \(-N_f\). Indeed, near that collar the orientation is \(\operatorname{or}(M)\,ds\,dt\); contraction with the outward normal \(-\partial_s\) gives \(\operatorname{or}(M)\,dt\). At \(t=0,1\) the other boundary pieces are \(-K\) and \(+K\), respectively, with exactly the displayed identification. The class \(x\) extends, as does the preferred spin structure (choose the appropriate one of the two spin structures of the bordism). This verifies the direction and the orientation of the comparison, not only equality of the unoriented manifolds.

\item It follows that \[
 x^2=k(f)y,\qquad p_1(TN_f)=\ell(f)y.                 \tag{P3F.22.2}
\]

\subsubsection{2. Removal of the cubic coordinate by an extending map}\label{removal-of-the-cubic-coordinate-by-an-extending-map}

\item Set \(f_0=f\circ a^{-k(f)}\). Jin's additive integral coordinates give \[
 k(f_0)=0,\qquad \ell(f_0)=\ell(f)-4k(f)=-24A(f).
\] The explicit quaternionic map in P3F.4 extends every lift of \(a^n\) over \(D^4_u\times S^3_y\). Apply that extension on the first gluing piece and the identity on the second. With the convention above it identifies \(\lambda(f_0)\) and \(\lambda(f)\) by an orientation-preserving diffeomorphism. This step does not assume in advance that \(\lambda\) is a group homomorphism.

\item We can therefore suppose \(k=0\). Write \(p=\ell(f_0)\), so that \(x^2=0\) and \(p_1(TN_{f_0})=py\).

\subsubsection{3. The surgery and its bounding eight-manifold}\label{the-surgery-and-its-bounding-eight-manifold}

\item Let \(L\to N_{f_0}\) have first Chern class \(x\), let \(E=S(L)\), and let \(D(L)\) be its oriented disc bundle. The preimage of the first \(K=S^2\times D^4\) is \(S^3\times D^4\). Perform the framed surgery which replaces that piece by \(D^4\times S^3\). Its framing is the product framing of the \(D^4\) factor, so the attaching map is completely specified. The result is a homotopy sphere, either by its explicit identification below or by the elementary Mayer--Vietoris computation in P3F.4.

\item The orientation is relevant. The oriented circle bundle decomposes as \[
 E=(S^3\times D^4)_+\cup_{\widetilde f_0}
                                      -(S^3\times D^4)_-.
\] The boundary of \(S^3\times D^4\) has orientation \(-\operatorname{or}(S^3\times S^3)\), whereas the boundary of \(D^4\times S^3\) has the positive product orientation. Thus the replacement has orientation \(-(D^4\times S^3)\). The oriented surgery sphere is exactly \(-\lambda(f_0)\).

\item An eight-manifold bounding that surgery sphere is \[
 W=D(L)\cup (D^4\times D^4),
\] where the four-handle attaches along the specified \(S^3\times D^4\). Equivalently it is obtained by replacing the Hopf-disc neighbourhood of \(S^2\) by \(\mathbb CP^2\times D^4\). This is the explicit construction in Jiang14, proof of Lemma3.5, Case2. It has an embedded \(\mathbb CP^2\) with trivial normal four-plane and a second four-cycle coming from the generator of \(H_4(N_{f_0})\). Their intersection matrix is \[
 \begin{pmatrix}0&1\\1&0\end{pmatrix}.              \tag{P3F.22.3}
\] The first self-intersection is zero by the trivial normal bundle; the second is zero because the disc-bundle Thom map \(H^2(N_{f_0})\to H^4(N_{f_0})\) is multiplication by \(x\), and \(x^2=0\). The mixed intersection is one because the core \(S^2\) is dual to \(y\). These facts also give a direct check of the intersection matrix in the cited proof.

\item Let \(z\in H^2(W;\mathbb Z)\) be the class extending \(x\), positive on the projective line in that \(\mathbb CP^2\). It lifts \(w_2(TW)\): on \(D(L)\) the latter is \(x\bmod2\), and on \(\mathbb CP^2\times D^4\) it is the hyperplane class modulo two. Let \(u,v\) be the relative degree-four classes dual to the two cycles in \(P3F.22.3\). The restrictions to those two cycles give \[
 z^2=jv,\qquad p_1(TW)=p\,ju+3jv,
\] where \(j:H^4(W,\partial W)\to H^4(W)\) is the natural map. Indeed the restrictions of \(z^2\) are \((1,0)\), and those of \(p_1\) are \((3,p)\). Hence \[
 \sigma(W)=0,\quad
 \langle p_1^2,[W,\partial W]\rangle=6p,\quad
 \langle z^2p_1,[W,\partial W]\rangle=p,\quad
 \langle z^4,[W,\partial W]\rangle=0.              \tag{P3F.22.4}
\]

\subsubsection{4. The invariant and the resulting homomorphism}\label{the-invariant-and-the-resulting-homomorphism}

\item For a spin boundary homotopy seven-sphere and this characteristic extension \(z\), the spin-c Eells--Kuiper formula is \[
 \mu(\partial W)=
 \frac{p_1^2-4\sigma(W)}{896}
 -\frac{z^2p_1}{192}+\frac{z^4}{384}\pmod{\mathbb Z}.
\] The formula is Jiang14 Section3.1\(v\), with its stated hypotheses; here \(\partial W\) is a homotopy sphere, so the required rational relative-to-absolute degree-four map is an isomorphism. Its normalization also follows by subtracting the spin-c index \(\int e^{z/2}\widehat A\) from the signature polynomial. Substitution of \(P3F.22.4\) gives \[
 \mu(-\lambda(f_0))=
 \frac{6p}{896}-\frac{p}{192}=\frac{p}{672}.
\] Since \(p=-24A(f)\), reversal of the boundary orientation proves \(P3F.22.1\).

\item The Eells--Kuiper invariant is additive under connected sum and identifies \(\Theta _7\) with \(\frac1{28}\mathbb Z/\mathbb Z\). For the positive normalization one can also check the Milnor sphere \(S_{2,-1}(\mathbb H)\): its disc bundle has Euler number one, Pontryagin number six and signature one, giving \((36-4)/896=1/28\). Davis82 Section3 computes the same value. Thus (P3F.22.1), together with Jin's additivity of \(A\), proves that \(\lambda\) is a homomorphism and gives its exact image and kernel. No classification statement about circle actions is needed for this last deduction.

\end{enumerate}
\end{proof}

% Original attached proof of night:P3F.24; retained solely as audit evidence.
\begin{proof}
\begin{enumerate}[label=\arabic*.]
\item Let \(\Phi\) be an orientation-preserving diffeomorphism of \(S^6\). Write \(\Sigma_\Phi=D^7\cup_\Phi(-D^7)\), with the direction of the boundary identification fixed from the first disc to the second. Its forward mapping torus satisfies \[
 T^+_\Phi=(S^6\times[0,1])/((z,1)\sim(\Phi(z),0))
       \cong(S^1\times S^6)\mathbin{\#}\Sigma_\Phi.   \tag{P3F.24.1}
\] Its reverse mapping torus, the convention used for the Jin six-dimensional mapping torus, satisfies \[
 T^-_\Phi=(S^6\times[0,1])/((z,0)\sim(\Phi(z),1))
       \cong(S^1\times S^6)\mathbin{\#}(-\Sigma_\Phi).\tag{P3F.24.2}
\] These are orientation-preserving diffeomorphisms for the indicated product and gluing orientations. For both spin structures on each mapping torus, its Eells--Kuiper invariant is respectively \(+\mu(\Sigma_\Phi)\) and \(-\mu(\Sigma_\Phi)\). In particular a nonzero such invariant excludes even an unoriented diffeomorphism with \(S^1\times S^6\).

\item \textbf{Proof.} First isotope \(\Phi\) to be the identity on a small oriented disc. The needed isotopy of disc embeddings is explicit: shrink the parametrized image disc toward its centre, move its centre along a path, move its oriented frame along the oriented frame bundle, and expand to the chosen target disc. Compactness gives a common sufficiently small radius along the path. Isotopy extension carries this isotopy to the ambient sphere. The result is a diffeomorphism equal to the identity on the target disc, and so supported in the complementary disc \(D^6\). Mapping tori of isotopic maps are diffeomorphic, by applying the isotopy along the interval; the same is true for their boundary-clutched spheres by extension over a collar. Thus this modification loses no data.

\item In the standard product \(S^1\times S^6\), choose a seam of the circle and the seven-ball \(B=D^6\times[-\epsilon,\epsilon]\) crossing that seam; round its corners. Its upper face has the positive orientation of \(D^6\), because the latter has even dimension. Define a boundary diffeomorphism of \(B\) to equal \(\Phi|_{D^6}\) on this face and the identity on the other face and the cylindrical collar. It is smooth since \(\Phi\) is the identity near \(\partial D^6\). Regluing this ball into its complement changes precisely the upper-to-lower seam identification to \(z\mapsto\Phi(z)\), and hence produces \(T^+_\Phi\). The boundary diffeomorphism, as an oriented map of \(\partial B\), is the original disc-supported \(\Phi\): both are its extension by the identity to the complementary six-disc.

\item The same regluing is connected sum with \(\Sigma_\Phi\). To check the map, remove an interior seven-ball from the second disc in \(D^7\cup_\Phi(-D^7)\). That remaining second piece is a collar \(S^6\times[0,1]\). Inserting this collar in the ball removed from \(S^1\times S^6\) gives exactly the specified boundary identification \(\Phi\). This proves (P3F.24.1) with its gluing direction. Replacing \(\Phi\) by \(\Phi^{-1}\) gives the reverse mapping torus and the reversed clutched sphere, proving \(P3F.24.2\).

\item There are exactly two spin structures because \(H^1(S^1\times S^6;\mathbb Z/2)=\mathbb Z/2\), and connected sum with a homotopy sphere does not change this group. Both product spin structures extend over \(S^1\times D^7\). That explicit eight-manifold has zero signature and zero first Pontryagin square. A homotopy seven-sphere bounds a spin eight-manifold; for instance the usual parallelizable representative supplied by \(\Theta _7=bP_8\) suffices. The boundary connected sum of this bounding manifold with \(S^1\times D^7\) bounds (P3F.24.1) for either spin structure. Its middle intersection form and its Pontryagin square are direct sums, since the joining handle has index one. Hence \[
 \mu(T^+_\Phi)=\mu(\Sigma_\Phi)
\] for each spin structure. The rational relative-to-absolute degree-four map is an isomorphism here because the mapping torus has \(H^3(-;\mathbb Q)=H^4(-;\mathbb Q)=0\). Thus the usual spin Eells--Kuiper expression \((p_1^2-4\sigma)/896\) applies without a choice of rational lift. Its invariance follows by gluing two such bounds: on the resulting closed spin eight-manifold this expression is the integral \(\widehat A\)-genus. Orientation reversal changes its sign. Both values for the product are zero, so neither preserving nor reversing orientation can identify a nonzero example with that product.

\end{enumerate}
\end{proof}
